\documentclass[12pt,a4paper]{article}
\usepackage[utf8]{inputenc}
\usepackage[T1]{fontenc}
\usepackage[english]{babel}
\usepackage{geometry}
\usepackage{hyperref}
\usepackage{booktabs}
\usepackage{longtable}
\usepackage{array}
\usepackage{listings}
\usepackage{xcolor}
\usepackage{parskip}
\usepackage{fancyhdr}
\usepackage{amsmath}
\usepackage{tcolorbox}
\usepackage{enumitem}
\usepackage{ragged2e}
\tcbuselibrary{skins,breakable}

\geometry{margin=2.5cm}
\hypersetup{colorlinks=true, linkcolor=blue!60!black, urlcolor=blue!60!black}

\definecolor{codebg}{RGB}{245,245,245}
\definecolor{noteblue}{RGB}{220,235,250}
\definecolor{warnred}{RGB}{255,235,230}
\definecolor{importantbg}{RGB}{230,250,235}
\definecolor{commentgreen}{RGB}{0,128,0}
\definecolor{keywordblue}{RGB}{0,0,180}
\definecolor{stringbrown}{RGB}{160,60,0}

\lstset{
  backgroundcolor=\color{codebg},
  basicstyle=\ttfamily\small,
  keywordstyle=\color{keywordblue}\bfseries,
  commentstyle=\color{commentgreen}\itshape,
  stringstyle=\color{stringbrown},
  breaklines=true, breakatwhitespace=true,
  frame=single, rulecolor=\color{gray!40},
  xleftmargin=6pt, xrightmargin=6pt,
  showstringspaces=false, language=Python
}

\newcolumntype{P}[1]{>{\RaggedRight\arraybackslash}p{#1}}

\pagestyle{fancy}
\fancyhf{}
\rhead{\texttt{rientra\_import.py} (Utenti --- R2RML)}
\lhead{\leftmark}
\cfoot{\thepage}

\title{\textbf{Rientra@ Import Tool --- Strada B}\\[0.4em]
       \large W3C RDB2RDF via R2RML\\[0.4em]
       \normalsize Technical Documentation of\\
       \texttt{Import/Utenti/rientra\_import.py}}
\author{STIIMA-CNR}
\date{Version 2.0 --- May 2026}

\begin{document}
\begin{titlepage}\centering\maketitle\end{titlepage}
\tableofcontents
\newpage

\clearpage
\section{Overview}

\texttt{Import/Utenti/rientra\_import.py} is the \textbf{Strada~B} (W3C-compliant path) of the Rientra@ patient import pipeline. Unlike the simpler \texttt{Lavori} variant which hard-codes RDF/XML string templates, this script implements a full \textbf{W3C RDB2RDF} workflow using:

\begin{enumerate}[noitemsep]
  \item \textbf{R2RML mapping} auto-generated by introspecting the ontology with \texttt{rdflib}.
  \item \textbf{Morph-KGC} as the R2RML materialisation engine, which produces an in-memory \texttt{rdflib.Graph}.
  \item \textbf{Text injection}: the materialised triples are serialised to RDF/XML string blocks and injected before \texttt{</rdf:RDF>} without re-parsing the ontology.
\end{enumerate}

\begin{tcolorbox}[colback=noteblue!40, colframe=blue!60!black, title=Design Principle: Standard Compliance]
The R2RML standard (\url{https://www.w3.org/TR/r2rml/}) and Direct Mapping (\url{https://www.w3.org/TR/rdb-direct-mapping/}) govern how relational data is transformed into RDF. Morph-KGC handles this transformation automatically from a declarative \texttt{.ttl} mapping file, removing the need for manually coded RDF builders.
\end{tcolorbox}

\subsection{Two Modes of Operation}

\begin{longtable}{@{}P{4cm}P{10.5cm}@{}}
\toprule
\textbf{Mode} & \textbf{Command} \\
\midrule
\endhead
Generate mapping &
  \texttt{python rientra\_import.py --generate-mapping --ontology Rientra.rdf} \\
Run import &
  \texttt{python rientra\_import.py --ontology Rientra.rdf --dataset dataset\_import.sql} \\
Run import (custom output) &
  \texttt{python rientra\_import.py --ontology Rientra.rdf --dataset dataset.sql --output Rientra\_v2.rdf} \\
Run import (custom mapping) &
  \texttt{python rientra\_import.py --ontology Rientra.rdf --dataset dataset.sql --mapping custom.ttl} \\
\bottomrule
\end{longtable}

\subsection{Dependencies}

\begin{lstlisting}
pip install morph-kgc rdflib
\end{lstlisting}

\begin{longtable}{@{}P{4cm}P{10.5cm}@{}}
\toprule
\textbf{Package} & \textbf{Purpose} \\
\midrule
\endhead
\texttt{morph-kgc} & W3C R2RML materialisation engine \\
\texttt{rdflib} & Ontology parsing (for mapping generation) and in-memory graph handling \\
\texttt{sqlite3} & In-memory SQL execution for the dataset \\
\texttt{argparse} & CLI argument parsing \\
\texttt{xml.sax.saxutils} & XML-safe escaping of literal values \\
\texttt{tempfile} & Temporary on-disk SQLite file for Morph-KGC \\
\bottomrule
\end{longtable}

\clearpage
\section{Module-Level Constants}

\begin{longtable}{@{}P{4cm}P{10.5cm}@{}}
\toprule
\textbf{Constant} & \textbf{Value / Role} \\
\midrule
\endhead
\texttt{NS\_FOAF} & \texttt{Namespace("http://www.stiima.cnr.it/FOAF-excerpt\#")} \\
\texttt{NS\_HC} & \texttt{Namespace("http://www.stiima.cnr.it/RientraHC\#")} \\
\texttt{NS\_ICF} & \texttt{Namespace("http://www.stiima.cnr.it/ICF-exc-coreset\#")} \\
\texttt{NS\_PERSON} & \texttt{Namespace("http://www.stiima.cnr.it/Person-CommonBox\#")} \\
\texttt{NS\_RIE} & \texttt{Namespace("http://www.stiima.cnr.it/RientraOnt3Merged\#")} \\
\texttt{CLOSING\_TAG} & \texttt{"</rdf:RDF>"} --- injection anchor in the RDF file \\
\texttt{QUALIFIER\_MAP} & \texttt{\{"b": "BFqual", "d": "AP1qual", "s": "BS1qual"\}} \\
\texttt{DEFAULT\_MAPPING} & \texttt{Path(\_\_file\_\_).parent / "rientra\_mapping.ttl"} \\
\bottomrule
\end{longtable}

Note: namespaces are \texttt{rdflib.Namespace} objects (not plain strings), unlike the \texttt{Lavori} variant.

\clearpage
\section{Part A --- R2RML Mapping Generator}

\subsection{Overview of \texttt{generate\_mapping}}

\textbf{Signature:} \texttt{generate\_mapping(ontology\_path: str, mapping\_out: str) -> None}

Introspects the OWL ontology using \texttt{rdflib} to discover all data properties declared for \texttt{FOAF:Person} and \texttt{HC\_Descriptor}, then generates a standards-compliant \texttt{rientra\_mapping.ttl} file. This file needs to be regenerated only when the ontology structure changes.

\subsection{Step 1 --- Property Classification}

The generator scans all \texttt{owl:DatatypeProperty} individuals and classifies them into four buckets:

\begin{longtable}{@{}P{3.5cm}P{2.5cm}P{8.5cm}@{}}
\toprule
\textbf{Bucket} & \textbf{Range} & \textbf{R2RML Mapping} \\
\midrule
\endhead
\texttt{plain\_props} & \texttt{xsd:string} or \texttt{xsd:anyURI} or None & \texttt{rr:objectMap [ rr:column "..." ]} \\
\texttt{typed\_props} & Any other XSD type & \texttt{rr:objectMap [ rr:column "...", rr:datatype xsd:... ]} \\
\texttt{boolean\_props} & \texttt{xsd:boolean} & \texttt{rr:objectMap [ rr:constant "false"\^{}\^{}xsd:boolean ]} \\
\texttt{qualifier\_props} & \texttt{HC\_Descriptor} domain & Routed to per-prefix views (\texttt{v\_desc\_b/d/s}) \\
\bottomrule
\end{longtable}

Only properties whose \texttt{rdfs:domain} is \texttt{FOAF:Person} are included in the Person map. Properties with \texttt{HC\_Descriptor} domain populate qualifier maps. \texttt{isSelected} (boolean) is always hardcoded to \texttt{false}.

\subsection{Step 2 --- Generated TriplesMap Structure}

The \texttt{.ttl} file contains \textbf{4 TriplesMap blocks}:

\begin{longtable}{@{}P{2cm}P{3.5cm}P{9cm}@{}}
\toprule
\textbf{Map ID} & \textbf{Source Table} & \textbf{Triples Produced} \\
\midrule
\endhead
\texttt{\#PersonMap} & \texttt{person} & \texttt{FOAF:Person} individual with all data properties + \texttt{isInHealthCondition} link \\
\texttt{\#HealthConditionMap} & \texttt{person} & \texttt{Health\_Condition} individual, one per person \\
\texttt{\#DescriptorMap\{B|D|S\}} & \texttt{v\_desc\_b/d/s} & One \texttt{HC\_Descriptor} per (person, ICF code) with qualifier \\
\texttt{\#DescriptorLinkMap} & \texttt{v\_hc\_links} & \texttt{isDescribedBy} links from HC to each descriptor \\
\bottomrule
\end{longtable}

\subsubsection*{IRI Templates Used}
\begin{longtable}{@{}P{5cm}P{9.5cm}@{}}
\toprule
\textbf{Entity} & \textbf{IRI Template} \\
\midrule
\endhead
Person & \texttt{http://www.stiima.cnr.it/Person-CommonBox\#\{person\_id\}} \\
Health Condition & \texttt{http://www.stiima.cnr.it/Person-CommonBox\#HC\{person\_id\}} \\
HC Descriptor & \texttt{http://www.stiima.cnr.it/RientraHC\#des\_\{person\_id\}\_\{icf\_code\}} \\
ICF Individual & \texttt{http://www.stiima.cnr.it/ICF-exc-coreset\#\{icf\_code\}} \\
\bottomrule
\end{longtable}

\clearpage
\section{Part B --- Import Pipeline Functions}

\subsection{\texttt{sanitize\_id(name) -> str}}
Replaces all characters that are not alphanumeric, underscore, or hyphen with \texttt{\_}. Ensures all generated IRIs are syntactically valid.

\subsection{\texttt{load\_dataset(sql\_path) -> sqlite3.Connection}}
Reads the SQL file as plain text and executes it via \texttt{conn.executescript(sql)} into an \textbf{in-memory SQLite database}. Returns the connection with \texttt{row\_factory = sqlite3.Row}. Exits with code~1 on SQL errors.

\subsection{\texttt{create\_views(conn, known\_icf) -> tuple[int, int]}}

Creates the auxiliary views required by the R2RML mapping:

\begin{enumerate}[noitemsep]
  \item Populates the temporary \texttt{\_known\_icf} table from the \texttt{known\_icf} set.
  \item Creates filtered views \texttt{v\_desc\_b}, \texttt{v\_desc\_d}, \texttt{v\_desc\_s} --- each selects rows from \texttt{hc\_descriptor} where the ICF code prefix matches and the code is present in \texttt{\_known\_icf}.
  \item Creates a \texttt{v\_hc\_links} view as a \texttt{UNION ALL} of the three prefix views.
\end{enumerate}

Returns \texttt{(valid\_count, skipped\_count)}.

\subsection{\texttt{extract\_known\_icf(rdf\_text) -> set[str]}}

Scans the raw RDF text with a regular expression to extract all ICF individual local names from the \texttt{ICF-exc-coreset} namespace. Used to filter out ICF codes not present in the ontology before running Morph-KGC.

\subsection{\texttt{person\_exists(rdf\_text, person\_id) -> bool}}

Checks whether the Person-CommonBox IRI for \texttt{person\_id} is already present in the RDF file text (substring search). Prevents duplicate injection.

\subsection{\texttt{export\_db\_to\_file(conn) -> str}}

Copies the in-memory SQLite database to a temporary file on disk using \texttt{conn.iterdump()} and returns the temp file path. Required because Morph-KGC needs a file-based SQLite URL (\texttt{sqlite:///...}).

\subsection{\texttt{run\_r2rml(mapping\_path, db\_file) -> Graph}}

Builds a Morph-KGC configuration string and calls \texttt{morph\_kgc.materialize(config)}, which returns an \texttt{rdflib.Graph} containing all materialised triples. The configuration specifies \texttt{N-TRIPLES} as output format and points to the SQLite file.

\begin{lstlisting}
config = """
[CONFIGURATION]
output_format=N-TRIPLES

[RientraDataSource]
mappings={mapping_path}
db_url=sqlite:///{db_file}
"""
return morph_kgc.materialize(config)
\end{lstlisting}

\subsection{\texttt{triples\_to\_rdfxml\_blocks(new\_graph, persons\_to\_add) -> list[str]}}

Converts the materialised \texttt{rdflib.Graph} into a list of RDF/XML string blocks ready for injection. For each new person:

\begin{enumerate}[noitemsep]
  \item \textbf{HC\_Descriptor blocks}: iterates \texttt{HC isDescribedBy Descriptor} triples; for each descriptor retrieves \texttt{involvesICFCode} and the appropriate qualifier (\texttt{BFqual}/\texttt{AP1qual}/\texttt{BS1qual}) from the graph.
  \item \textbf{Health\_Condition block}: one \texttt{owl:NamedIndividual} of type \texttt{HC:Health\_Condition} with all \texttt{isDescribedBy} links.
  \item \textbf{Person block}: one \texttt{owl:NamedIndividual} of type \texttt{FOAF:Person} with all plain and typed FOAF properties (XML-escaped), plus hardcoded \texttt{isSelected=false}.
\end{enumerate}

\clearpage
\section{Entry Point --- \texttt{main()}}

\subsection{Command-Line Arguments}

\begin{longtable}{@{}P{4.5cm}P{2cm}P{8cm}@{}}
\toprule
\textbf{Argument} & \textbf{Required} & \textbf{Description} \\
\midrule
\endhead
\texttt{--ontology} & Yes & Path to the \texttt{.rdf} ontology file \\
\texttt{--generate-mapping} & No & If set, runs \texttt{generate\_mapping()} and exits \\
\texttt{--dataset} & Conditional & Path to the \texttt{.sql} dataset (required unless \texttt{--generate-mapping}) \\
\texttt{--output} & No & Output path; defaults to overwriting \texttt{--ontology} \\
\texttt{--mapping} & No & Path to the \texttt{.ttl} mapping file; defaults to \texttt{rientra\_mapping.ttl} \\
\bottomrule
\end{longtable}

\subsection{Six-Step Import Execution}

\begin{longtable}{@{}P{2cm}P{12.5cm}@{}}
\toprule
\textbf{Step} & \textbf{Description} \\
\midrule
\endhead
1/6 & Load SQL dataset into in-memory SQLite; count persons. \\
2/6 & Read ontology as raw text; extract known ICF codes; identify new vs already-present persons. \\
3/6 & Create auxiliary views (\texttt{v\_desc\_b/d/s}, \texttt{v\_hc\_links}) filtered by \texttt{\_known\_icf}. \\
4/6 & Export in-memory DB to temp file; run Morph-KGC R2RML materialisation. \\
5/6 & Convert materialised graph to RDF/XML blocks; print per-person descriptor count. \\
6/6 & Create backup (\texttt{.bak}); inject blocks before \texttt{</rdf:RDF>}; write output file. \\
\bottomrule
\end{longtable}

\subsection{Injection Mechanism}

\begin{lstlisting}
injection = (
    "\n\n\n    <!-- ===== PAZIENTI IMPORTATI (RDB2RDF via R2RML) ===== -->\n\n"
    + "\n\n".join(blocks)
    + "\n\n"
)
updated = rdf_text.replace(CLOSING_TAG, injection + CLOSING_TAG, 1)
Path(out_path).write_text(updated, encoding="utf-8")
\end{lstlisting}

The \texttt{str.replace(..., 1)} call replaces only the first occurrence of \texttt{</rdf:RDF>}, which is always the last tag in a well-formed RDF/XML file.

\clearpage
\section{Function Reference Summary}

\begin{longtable}{@{}P{6.5cm}P{3cm}P{5.5cm}@{}}
\toprule
\textbf{Function} & \textbf{Returns} & \textbf{Part} \\
\midrule
\endhead
\texttt{generate\_mapping(ontology, out)} & \texttt{None} & A --- Mapping Generator \\
\texttt{sanitize\_id(name)} & \texttt{str} & B --- Import Pipeline \\
\texttt{load\_dataset(sql\_path)} & \texttt{Connection} & B --- Import Pipeline \\
\texttt{create\_views(conn, known\_icf)} & \texttt{tuple[int,int]} & B --- Import Pipeline \\
\texttt{extract\_known\_icf(rdf\_text)} & \texttt{set[str]} & B --- Import Pipeline \\
\texttt{person\_exists(rdf\_text, pid)} & \texttt{bool} & B --- Import Pipeline \\
\texttt{export\_db\_to\_file(conn)} & \texttt{str} & B --- Import Pipeline \\
\texttt{run\_r2rml(mapping, db\_file)} & \texttt{Graph} & B --- Import Pipeline \\
\texttt{triples\_to\_rdfxml\_blocks(g, pids)} & \texttt{list[str]} & B --- Import Pipeline \\
\texttt{main()} & \texttt{None} & Entry point \\
\bottomrule
\end{longtable}

\clearpage
\section{Comparison with Lavori Variant}

\begin{longtable}{@{}P{4cm}P{5.5cm}P{5.5cm}@{}}
\toprule
\textbf{Aspect} & \textbf{Utenti (Strada B)} & \textbf{Lavori (Strada A)} \\
\midrule
\endhead
Standard & W3C R2RML + Direct Mapping & None (ad-hoc) \\
Mapping & Auto-generated \texttt{.ttl} from ontology & Hard-coded Python string templates \\
Materialisation & Morph-KGC (external engine) & Manual \texttt{build\_*\_xml()} functions \\
Namespaces & \texttt{rdflib.Namespace} objects & Plain strings \\
Property discovery & Introspective (reads ontology) & Hard-coded \texttt{FOAF\_DATATYPES} dict \\
ICF filtering & SQL VIEW + \texttt{\_known\_icf} table & \texttt{iri\_exists()} substring check \\
S-prefix ICF codes & Supported (\texttt{BS1qual}) & Not present \\
Injection mechanism & Same (string replace before \texttt{</rdf:RDF>}) & Same \\
Dataset format & \texttt{.sql} file (SQLite executescript) & \texttt{.sql} file (SQLite executescript) \\
Lines of code & 578 & 258 \\
\bottomrule
\end{longtable}

\clearpage
\section{Known Issues and Notes}

\begin{tcolorbox}[colback=noteblue!40, colframe=blue!60!black,
                  title=Note: Mapping must be regenerated after ontology changes]
The \texttt{rientra\_mapping.ttl} file is generated once from the ontology structure. If new data properties are added to \texttt{FOAF:Person} or \texttt{HC\_Descriptor} in \texttt{Rientra.rdf}, the mapping must be regenerated with \texttt{--generate-mapping} before the next import. The generated file includes a comment block with the exact command to re-run.
\end{tcolorbox}

\vspace{0.6em}

\begin{tcolorbox}[colback=warnred!40, colframe=red!60!black,
                  title=Warning: Morph-KGC requires a file-based SQLite database]
\texttt{morph\_kgc.materialize()} cannot connect to an in-memory SQLite database (\texttt{:memory:}). The function \texttt{export\_db\_to\_file()} copies the in-memory DB to a \texttt{tempfile.NamedTemporaryFile}. The temp file is deleted in the \texttt{finally} block via \texttt{Path(db\_file).unlink(missing\_ok=True)}.
\end{tcolorbox}

\vspace{0.6em}

\begin{tcolorbox}[colback=importantbg!60, colframe=green!50!black,
                  title=Important: E-prefix ICF codes are not supported]
Environmental Factor codes (\texttt{e}-prefix) are not imported. In the \texttt{Lavori} variant this is enforced by explicit \texttt{prefix == "e"} check; in the \texttt{Utenti} variant it is enforced implicitly because \texttt{QUALIFIER\_MAP} only contains \texttt{b}, \texttt{d}, \texttt{s} prefixes, and \texttt{create\_views} creates views only for those three prefixes. Rows with \texttt{e}-prefix codes are excluded from \texttt{\_known\_icf} join and counted as skipped.
\end{tcolorbox}

\end{document}
